\documentclass[11pt]{article}
\usepackage[margin=1in]{geometry}
\usepackage{amsmath,amssymb,amsthm}
\usepackage[T1]{fontenc}
\usepackage{lmodern}
\usepackage{hyperref}
\usepackage{enumitem}
\hypersetup{colorlinks=true,linkcolor=blue,urlcolor=blue}
\setlength{\emergencystretch}{2em}
\newtheorem{theorem}{Theorem}
\newtheorem{lemma}[theorem]{Lemma}
\newtheorem{proposition}[theorem]{Proposition}
\newcommand{\D}{\mathbb D}
\newcommand{\C}{\mathbb C}
\newcommand{\A}{A^2(\D)}
\newcommand{\ip}[2]{\langle #1,#2\rangle}
\newcommand{\norm}[1]{\lVert #1\rVert}
\newcommand{\Lean}[1]{\texttt{\detokenize{#1}}}
\title{The Bergman shift has exactly two reducing subspaces\\
\large A disproof of conjecture 00000000930}
\author{Independent mathematical submission}
\date{}
\begin{document}
\maketitle

\begin{abstract}
For the ordinary unweighted Bergman shift, the closed complex reducing
subspaces are exactly $\{0\}$ and the whole Bergman space. Their cardinality is
therefore $2$, not the cardinality of the continuum. We prove this directly on
square-integrable holomorphic functions by division by $z-w$. The accompanying
Lean 4 project also proves the analytic completeness bridge, the equivalence of
two standard reducing-subspace definitions, and the exact cardinality statement.
\end{abstract}

\section{Statement, conventions, and conclusion}
The supplied English statement is:
\begin{quote}
Definition: The Bergman shift ($M_z$ on $A^2$). Conjecture: The lattice of
reducing subspaces of the Bergman shift has the cardinality of the continuum,
each indexed by Blaschke products (a complete classification of reducing
subspaces).
\end{quote}
The Chinese statement specifies the same operator and the same
continuum-cardinality and Blaschke-indexing assertions.

Let $\D=\{z\in\C:|z|<1\}$ and let $dA=dx\,dy$ be unnormalized planar
Lebesgue area. Set
\[
 \A=\left\{ f:\D\to\C:\ f\text{ is holomorphic and }
                    \int_\D |f(z)|^2\,dA(z)<\infty\right\}.
\]
Our inner product is conjugate-linear in the first variable:
\[
 \ip f g=\int_\D \overline{f(z)}g(z)\,dA(z).
\]
Using $dA/\pi$ instead multiplies this inner product by a positive constant
and leaves its topology, orthogonality, and reducing subspaces unchanged.
The operator in question is exactly
\[
 S=M_z,\qquad (Sf)(z)=zf(z).
\]
It is bounded, with $\norm{Sf}_2\leq\norm f_2$.

A reducing subspace means a closed complex linear subspace $M\subseteq\A$
for which $S(M)\subseteq M$ and $S(M^\perp)\subseteq M^\perp$.
Equivalently, it is a closed complex linear subspace invariant under both
$S$ and its Hilbert-space adjoint $S^*$; this equivalence is proved below and
in Lean.

\begin{theorem}\label{thm:main}
The complete collection of closed complex reducing subspaces of $M_z$ on
$\A$ is $\{\{0\},\A\}$. In particular, its cardinality is exactly $2$, and
is not the cardinality $\mathfrak c$ of the continuum.
\end{theorem}
This disproves a necessary, explicit clause of the conjecture and hence
disproves the conjecture as a conjunction. The Blaschke-product sentence does
not specify an indexing map or an equivalence relation on labels. No such map
or equivalence relation is invented here; the disproof is the failure of the
continuum-cardinality clause for the actual operator.

\section{The analytic Hilbert-space bridge}
For completeness we give the analytic facts required to apply orthogonal
projection to every closed subspace in the statement. These facts are also
proved in the accompanying project, rather than supplied as hypotheses.

\begin{lemma}[Area mean value]\label{lem:mean}
If $f$ is holomorphic on a neighborhood of $\overline{B(w,r)}$ and $r>0$, then
\[
 \int_{B(w,r)}f(z)\,dA(z)=\pi r^2 f(w).
\]
\end{lemma}
\begin{proof}
Cauchy's integral formula, with $z=w+\rho e^{i\theta}$, gives
\[
 \int_0^{2\pi} f(w+\rho e^{i\theta})\,d\theta=2\pi f(w)
 \quad(0<\rho\leq r).
\]
Integrating radially and using polar coordinates and Fubini gives
\[
 \int_{B(w,r)}f\,dA
 =\int_0^r\rho\!\int_{-\pi}^{\pi}f(w+\rho e^{i\theta})\,d\theta\,d\rho
 =\pi r^2 f(w).
\]
Changing the angular interval does not change the integral because the
integrand is $2\pi$-periodic. Continuity on the closed disk makes the polar
integrand integrable on its compact parameter rectangle. Translation handles
an arbitrary center $w$.
\end{proof}

\begin{lemma}[Evaluation estimate]\label{lem:eval}
Write $C=\norm{1}_{L^2(\D,dA)}$, a fixed finite constant. If
$\overline{B(w,r)}\subset\D$ and $r>0$, then every $f\in\A$ satisfies
\[
 |f(w)|\leq\frac{C}{\pi r^2}\norm f_2.
\]
\end{lemma}
\begin{proof}
The area mean-value formula, the triangle inequality, and Cauchy--Schwarz give
\[
 \pi r^2|f(w)|
 \leq\int_{B(w,r)}|f|\,dA
 \leq\int_\D|f|\,dA
 \leq C\norm f_2.
\]
The final inequality is the real $L^2$ inner-product inequality applied to
$1$ and $|f|$.
\end{proof}

\begin{proposition}\label{prop:hilbert}
The natural inclusion of $\A$ into $L^2(\D,dA)$ is injective and has closed
range. Consequently $\A$ is a Hilbert space with the displayed inner product.
\end{proposition}
\begin{proof}
Two continuous functions on the open disk that agree almost everywhere agree
everywhere: if their difference were nonzero at a point, continuity would
make it nonzero on a nonempty open set of positive area. This proves
injectivity.

Suppose $f_n\in\A$ converge to $u$ in $L^2$. They are $L^2$-Cauchy. Fix
$0<t<1$ and put $r=(1-t)/2>0$. For every $w\in B(0,t)$,
$\overline{B(w,r)}\subset\D$. Lemma~\ref{lem:eval} applied to $f_m-f_n$
therefore gives
\[
 \sup_{w\in B(0,t)}|f_m(w)-f_n(w)|
 \leq\frac{C}{\pi r^2}\norm{f_m-f_n}_2.
\]
Thus $(f_n)$ is uniformly Cauchy on every smaller disk. It has a pointwise
limit $g$ on $\D$, with locally uniform convergence. The locally uniform
limit theorem for holomorphic functions makes $g$ holomorphic.

Convergence in $L^2$ implies convergence in measure, so some subsequence
converges almost everywhere to a representative of $u$. The same subsequence
converges pointwise on $\D$ to $g$. Uniqueness of limits gives $g=u$ almost
everywhere. Hence $g$ is square integrable and represents $u$, proving that
the range is closed. Completeness follows from completeness of $L^2$.
\end{proof}

\section{Local division and the commutant}
\begin{lemma}[Division inside $A^2$]\label{lem:division}
For $f\in\A$ and $w\in\D$, define
\[
 q(z)=\begin{cases}
 \dfrac{f(z)-f(w)}{z-w},&z\ne w,\\[4pt]
 f'(w),&z=w.
 \end{cases}
\]
Then $q\in\A$ and
\[
 f-f(w)1=(S-wI)q.
\]
\end{lemma}
\begin{proof}
The removable singularity theorem makes $q$ holomorphic. Its continuity at
$w$ gives numbers $r>0$ and $C_0\geq0$ such that $|q(z)|\leq C_0$ for
$|z-w|<r$. Outside that neighborhood,
\[
 |q(z)|\leq r^{-1}\bigl(|f(z)|+|f(w)|\bigr).
\]
Thus everywhere on the disk,
\[
 |q(z)|\leq C_0+r^{-1}\bigl(|f(z)|+|f(w)|\bigr).
\]
The right side is in $L^2$, since $f\in L^2$ and the disk has finite area.
Therefore $q\in\A$. The stated identity holds pointwise, including at $w$.
\end{proof}

\begin{lemma}\label{lem:multiplier}
Let $P:\A\to\A$ be any complex-linear map satisfying $PS=SP$. Then
\[
 (Pf)(w)=f(w)(P1)(w)\qquad(f\in\A,\ w\in\D).
\]
\end{lemma}
\begin{proof}
Apply $P$ to the identity in Lemma~\ref{lem:division} and commute it with
$S-wI$:
\[
 Pf-f(w)P1=(S-wI)Pq.
\]
Evaluation at $w$ annihilates the right side.
\end{proof}
No boundedness assumption on $P$ is used in this lemma.

\begin{proposition}\label{prop:idempotent}
If $P:\A\to\A$ is complex linear, $PS=SP$, and $P^2=P$, then $P=0$ or
$P=I$.
\end{proposition}
\begin{proof}
Put $h=P1$. By Lemma~\ref{lem:multiplier}, $P$ acts pointwise as
multiplication by $h$. The equation $P^2 1=P1$ gives
\[
 h(w)(h(w)-1)=0\qquad(w\in\D).
\]
Both $h$ and $h-1$ are holomorphic, and the disk is connected. The identity
principle therefore implies that $h$ vanishes identically or that $h-1$
vanishes identically. The pointwise formula for $Pf$ then gives $P=0$ or
$P=I$, respectively.
\end{proof}

\section{Every closed reducing subspace and the exact cardinality}
First verify the equivalence of reducing-subspace conventions. If $M$ is
closed, Proposition~\ref{prop:hilbert} gives $M^{\perp\perp}=M$. If
$S(M^\perp)\subseteq M^\perp$, then for $x\in M$ and $y\in M^\perp$,
\[
 \ip{y}{S^*x}=\ip{Sy}{x}=0,
\]
so $S^*x\in M^{\perp\perp}=M$. Conversely, if $S^*(M)\subseteq M$, then
for $x\in M^\perp$ and $y\in M$,
\[
 \ip{y}{Sx}=\ip{S^*y}{x}=0,
\]
so $Sx\in M^\perp$. Together with $S(M)\subseteq M$, these are exactly
the two stated definitions.

\begin{proof}[Proof of Theorem~\ref{thm:main}]
Let $M$ be any closed reducing subspace. The Hilbert projection theorem
provides the orthogonal direct sum
\[
 \A=M\oplus M^\perp.
\]
Let $P$ be projection onto the first summand. For $u\in M$ and
$v\in M^\perp$, the two invariances give
\[
 PS(u+v)=Su=SP(u+v).
\]
Thus $P$ is a complex-linear idempotent commuting with $S$.
Proposition~\ref{prop:idempotent} implies $P=0$ or $P=I$, hence
$M=\{0\}$ or $M=\A$.

Both subspaces are closed and reducing. They are distinct because the
constant function $1$ belongs to $\A$ and is not zero. Sending
$\mathrm{false}$ to $\{0\}$ and $\mathrm{true}$ to $\A$ is consequently a
bijection from the two-element Boolean type onto the entire collection of
reducing subspaces. The cardinality is exactly $2$. Since every finite
cardinal is strictly smaller than $\mathfrak c$, it is not $\mathfrak c$.
\end{proof}

\section{Formalization and reproducibility}
The portable project is in \Lean{lean/}. It pins Lean 4.19.0 and
Mathlib v4.19.0, at commit
\begin{center}
\small\Lean{c44e0c8ee63ca166450922a373c7409c5d26b00b}
\end{center}
Run \Lean{lake build} from that directory. The
project-level option \Lean{warningAsError = true} makes warnings fatal.
All included project modules are reachable from the default target.

The definitions preserve the original analytic objects:
\begin{itemize}[leftmargin=*]
\item \Lean{Bergman930.space} consists of actual complex functions,
  holomorphic on \Lean{disk} and in \Lean{MemLp} at exponent $2$ for
  restricted planar Lebesgue measure. They are zero outside the disk solely
  to choose a unique representative of a function defined on $\D$.
\item \Lean{toL2_injective} and \Lean{functionEquivL2} identify those
  functions exactly with \Lean{l2Space}, the range in genuine area-$L^2$.
  \Lean{inner_toL2} proves the displayed integral inner-product formula.
  \Lean{l2Space_isClosed} proves closedness and supplies its completeness.
\item \Lean{shift} is pointwise multiplication by $z$;
  \Lean{bergmanShift} is this same operator transported to \Lean{l2Space},
  with its contraction bound checked. The reducing predicate uses this
  space's inherited $L^2$ topology and orthogonal complement.
\end{itemize}

\begin{samepage}
The principal checked statements, all in namespace \Lean{Bergman930}, are:
\begin{center}
\begin{tabular}{ll}
Mathematical assertion & Lean declaration\\\hline
Local division & \Lean{division_identity}\\
Commuting maps act pointwise & \Lean{commuting_map_apply}\\
Commuting idempotents are trivial & \Lean{commuting_idempotent_eq_zero_or_id}\\
Area mean value & \Lean{disk_mean}\\
$L^2$ evaluation estimate & \Lean{evaluation_bound}\\
Actual Bergman completeness & \Lean{l2Space_isClosed}\\
Adjoint-definition equivalence & \Lean{isReducing_iff_adjoint}\\
All closed reducing subspaces & \Lean{reducing_eq_bot_or_top}\\
Exact cardinality $2$ & \Lean{reducing_cardinality}\\
Failure of continuum cardinality & \Lean{not_continuum}
\end{tabular}
\end{center}
\end{samepage}

In particular, the final theorem is
\[
 \Lean{Cardinal.mk Bergman930.ReducingSubspace}\ne
 \Lean{Cardinal.continuum}.
\]
The preceding exact-cardinality theorem has right side $2$.
The type \Lean{ReducingSubspace} is a subtype of \emph{all} complex
submodules of the genuine Bergman Hilbert space satisfying
\Lean{IsReducing}; it is not an enumeration used as a definition.

\Lean{Verification.lean} prints the concrete definitions, the principal
theorem types, and their axiom dependencies. Fresh compilation of all owned
source modules succeeds with warnings treated as errors. The inspected
statements depend only on the standard Lean axioms
\Lean{propext}, \Lean{Classical.choice}, and \Lean{Quot.sound}.
There are no custom axioms, unfinished proofs, \Lean{native_decide} calls,
or unsafe or partial proof dependencies in the submitted source.

The trust boundary is explicit: the supplied precompiled objects of the
pinned standard dependencies are trusted imports; every owned source module
is freshly elaborated and checked. Three standard Mathlib modules absent
from that local precompiled cache were also freshly compiled from unchanged
pinned source. The final portable project uses their canonical Mathlib
imports and includes no renamed copies or dependency cache. Exact commands,
exit codes, raw outputs, and source hashes accompany the verification
record. No numerical computation is needed for the disproof.

\end{document}
